%% LyX 2.3.6 created this file.  For more info, see http://www.lyx.org/.
%% Do not edit unless you really know what you are doing.
\documentclass[english,aspectratio=169]{beamer}
\usepackage{lmodern}
\renewcommand{\sfdefault}{lmss}
\renewcommand{\ttdefault}{lmtt}
\usepackage[T1]{fontenc}
\usepackage[latin9]{inputenc}
\setlength{\parskip}{\medskipamount}
\setlength{\parindent}{0pt}
\usepackage{babel}
\usepackage{textcomp}
\usepackage{amssymb}
\usepackage{graphicx}
\PassOptionsToPackage{normalem}{ulem}
\usepackage{ulem}
\ifx\hypersetup\undefined
  \AtBeginDocument{%
    \hypersetup{unicode=true}
  }
\else
  \hypersetup{unicode=true}
\fi

\makeatletter

%%%%%%%%%%%%%%%%%%%%%%%%%%%%%% LyX specific LaTeX commands.
\pdfpageheight\paperheight
\pdfpagewidth\paperwidth

%% Because html converters don't know tabularnewline
\providecommand{\tabularnewline}{\\}

%%%%%%%%%%%%%%%%%%%%%%%%%%%%%% Textclass specific LaTeX commands.
% this default might be overridden by plain title style
\newcommand\makebeamertitle{\frame{\maketitle}}%
% (ERT) argument for the TOC
\AtBeginDocument{%
  \let\origtableofcontents=\tableofcontents
  \def\tableofcontents{\@ifnextchar[{\origtableofcontents}{\gobbletableofcontents}}
  \def\gobbletableofcontents#1{\origtableofcontents}
}

%%%%%%%%%%%%%%%%%%%%%%%%%%%%%% User specified LaTeX commands.
%\usetheme{CambridgeUS}
%\usetheme{Singapore}
\usetheme{Warsaw}
\usecolortheme{beaver}
\hypersetup{pdfpagemode=None}
\usepackage{tikz}
\usepackage{color}
\usepackage{listings}
\usepackage{multicol}
\setbeamertemplate{footline}[frame number]

\makeatother

\begin{document}
\title[ESS methods]{Earth energy budget}
\author{Magata J Mangatane}
\institute[UCT]{Department of Oceanography\\University of Cape Town\\mngmag002@myuct.ac.za}
\date{SEA3004F}
\makebeamertitle

\section{Earth energy budget}

\begin{frame}{Where does the energy go?}
\begin{columns}[t]

\column{8cm}
\begin{itemize}
\item <1->{\footnotesize{}First Law of Thermodynamics states that energy cannot be created or destroyed, only transformed or transferred.}
\item <2->{\footnotesize{}The change in internal energy of a system equals the net heat transfer into the system minus the net work done by the system}
\item <3->{\footnotesize{}In equation form: }\textbf{\footnotesize{}$\Delta U = Q - W$}{\footnotesize{}.}
\item <4->{\footnotesize{}This formulation is a specific expression of the conservation of energy}{\footnotesize\par}
{\footnotesize{}\uline{\href{https://teksguide.org/resource/151-first-law-thermodynamics}{Further Reading}}}{\footnotesize{}.}{\footnotesize\par}
\end{itemize}

\column{7cm}

\includegraphics[width=7cm]{SEA3004/2026/figures/Energy_conservation}
\end{columns}

\end{frame}
%
\begin{frame}{Recap: the global heat budget}

\includegraphics[height=0.75\paperheight]{SEA3004/2026/figures/energy_budget}
\end{frame}

\begin{frame}{Where does the energy go?}
\begin{columns}[t]

\column{8cm}
\begin{itemize}
\item <1->{\footnotesize{}Human-induced atmospheric composition changes
cause a radiative imbalance at the top of the Earth atmosphere which
is driving global warming (}{\footnotesize{}\uline{\href{https://doi.org/10.5194/essd-12-2013-2020}{von Schuckmann et al., 2020}}}{\footnotesize{}).
This quantity is called the }\textbf{\footnotesize{}Earth Energy Imbalance
(EEI)}{\footnotesize{}, and it is largely driven by the radiative
effect of excessive CO$_{2}$.} 
\item <2->{\footnotesize{}A zero-balance would maintain a stable }\textbf{\footnotesize{}planetary
temperature}{\footnotesize{}. A surplus leads to heat storage in various
parts of the ES and warming.} 
\item <3->{\footnotesize{}Knowing this imbalance would inform us how much
CO$_{2}$ we need to remove to attain the target temperature of the
}{\footnotesize{}\uline{\href{https://unfccc.int/process-and-meetings/the-paris-agreement/the-paris-agreement}{Paris Agreement}}}{\footnotesize{}.}{\footnotesize\par}
\end{itemize}

\column{7cm}

\includegraphics[width=7cm]{SEA2004/figures/EEI_fig8}
\end{columns}

\end{frame}
%


\begin{frame}{The radiative balance at the top of the atmosphere (TOA)}

\begin{columns}[t]


\column{7cm}
\begin{itemize}
\item <1->{\footnotesize{}(a) The red solid line is the intensity of incoming
solar radiation and the red dashed line is the intensity of radiation
lost to space (both determined using satellite-borne radiometer).
(i) the northern summer and (ii) the southern summer}{\footnotesize\par}
\item <2->{\footnotesize{}(b) The average temperature of surface waters
at different latitudes. At a given latitude, there will be surface
waters whose annual mean temperatures are higher or lower than shown
by the curve; this range is represented by the thickness of the pink
envelope. }\emph{\footnotesize{}(think about the reason why the x-axis
is not linear...)}{\footnotesize\par}
\end{itemize}

\column{5cm}

\includegraphics[width=5cm]{SEA2004/figures/radiation_balance}
\end{columns}

\end{frame}

\begin{frame}{OLR: Outgoing Longwave Radiation}

\begin{columns}[t]


\column{6cm}
\begin{itemize}
\item {\scriptsize{}The OLR is the total IR at wavelengths of 5 to 100 $\mu$m
that escapes from the top of Earth\textquoteright s atmosphere back
into space.}{\scriptsize\par}
\item {\scriptsize{}Assuming blackbody radiation, the OLR is a measure of
the temperature of the emitting bodies. }{\scriptsize\par}
\item {\scriptsize{}Measured OLR is not referred to the surface and must
be interpreted with due care. High OLR is usually due to warm surface
or low clouds (usually found on the Atlantic coast of SA). Low OLR
is due to the top of deep convective clouds which is cold (regions
of heavy rainfall). }{\scriptsize\par}
\end{itemize}

\column{7cm}
\begin{center}
\vspace{-1cm}
\includegraphics[width=6cm]{SEA2004/figures/OLR}
\par\end{center}

\end{columns}

\end{frame}

\begin{frame}{Surface heat flux balance over the ocean}

\begin{columns}[t]


\column{7.5cm}
\begin{itemize}
\item <1->The heat balance at the ocean surface is
\[
Q_{T}=Q_{s}+Q_{b}+Q_{e}+Q_{h}
\]

\begin{itemize}
\item <2->$Q_{s}$ = net rate of incoming (solar) shortwave radiation
\item <3->$Q_{b}$ = net rate of heat loss by the ocean as longwave radiation 
\item <4->$Q_{h}$ = rate of heat loss or gain due to conduction at the
surface (sensible heat flux)
\item <5->$Q_{e}$ = rate of heat loss or gain by the ocean through evaporation
or condensation (latent heat flux) 
\end{itemize}
\end{itemize}

\column{6cm}
\begin{center}
\includegraphics[width=6.5cm]{SEA2004/figures/heat_balance_surface}
\par\end{center}

\end{columns}

\end{frame}

\begin{frame}{Water heat and water temperature}

\begin{columns}[c]


\column{7.5cm}
\begin{itemize}
\item {\scriptsize{}Heat is a measure of how many molecules are vibrating
and how rapidly they are vibrating. Temperature records only how rapidly
the molecules of a substance are vibrating. }\textbf{\scriptsize{}Temperature
is an object\textquoteright s response to the input (or removal) of
heat}{\scriptsize\par}
\end{itemize}
\begin{center}
\vspace{-0.2cm}
\includegraphics[width=6.5cm]{SEA2004/figures/heat_capacity}
\par\end{center}

\column{7cm}

\begin{onlyenv}<2->
\begin{center}
{\tiny{}\vspace{-1cm}
}%
\begin{tabular}{|c|c|}
\hline 
\textbf{\tiny{}Substance} & \textbf{\tiny{}Heat capacity}\tabularnewline
 & {\tiny{}cal g$^{-1}$ $K^{-1}$}\tabularnewline
\hline 
\hline 
{\tiny{}Silver} & {\tiny{}0.06}\tabularnewline
\hline 
{\tiny{}Granite (or sand)} & {\tiny{}0.20}\tabularnewline
\hline 
{\tiny{}Aluminum} & {\tiny{}0.22}\tabularnewline
\hline 
{\tiny{}Alcohol (ethyl)} & {\tiny{}0.30}\tabularnewline
\hline 
{\tiny{}Gasoline} & {\tiny{}0.50}\tabularnewline
\hline 
\textbf{\tiny{}Pure water} & \textbf{\tiny{}1.00}\tabularnewline
\hline 
{\tiny{}Ammonia (liquid)} & {\tiny{}1.13}\tabularnewline
\hline 
\end{tabular}{\tiny\par}
\par\end{center}
\begin{itemize}
\item {\scriptsize{}Water has an extraordinarily high specific heat capacity
($c_{p}$, constant pressure): it resists changing temperature when
heat is added or removed. During phase changes, a large amount of
heat is required before seeing a change in temperature: these amounts
are the ``hidden'' }\textbf{\scriptsize{}latent heat of fusion and
vaporization}{\scriptsize\par}
\end{itemize}
\end{onlyenv}
\end{columns}

\end{frame}


\begin{frame}{Seasonal cycle of short wave radiation {[}W m$^{2}${]}}

\includegraphics[scale=0.45]{SEA2004/figures/shortwave}
\end{frame}

\begin{frame}{Zonally averaged annual fluxes {[}PW{]}}

\includegraphics[scale=0.45]{SEA2004/figures/zonal_fluxes}
\end{frame}

\begin{frame}{Surplus and deficit of energy}

\includegraphics[scale=0.3]{SEA2004/figures/heat_surplus}
\end{frame}


\subsection{Ocean heat transport}


\begin{frame}{Energy transfer is obtained through circulation}

\begin{itemize}
\item <1->The solar input drives the oceanic circulation through the atmosphere.
Both are affected by the Earth rotation. Energy (momentum) is transferred
from winds to the upper layers of the ocean through \textbf{frictional
coupling} between the ocean and the atmosphere at the sea surface
\item <2->The Sun also drives ocean circulation by causing variations in
the temperature and salinity of seawater which in turn control its
density. These processes drive the \textbf{thermohaline circulation}.
Changes in temperature are caused by fluxes of heat across the air-sea
boundary; changes in salinity are due to evaporation and precipitation,
but also, in polar regions, by the freezing and melting of ice. 
\item <3->All of these processes are linked directly or indirectly to the
effect of solar radiation and the related \textbf{heat and momentum
fluxes} at the surface. (suggested reading: Ocean Circulation, Open
University)
\end{itemize}
\end{frame}

\begin{frame}{Energy transport in the Earth system}

\begin{columns}[t]


\column{5cm}
\begin{itemize}
\item {\footnotesize{}The }\textbf{\footnotesize{}engine}{\footnotesize{}
of the Earth system is the }\textbf{\footnotesize{}Sun }{\footnotesize{}and
the Earth is in}\textbf{\footnotesize{} energy balance}{\footnotesize\par}
\item {\footnotesize{}The }\textbf{\footnotesize{}large scale circulation}{\footnotesize{}
of the ocean and atmosphere can be described by applying physical
principles of }\textbf{\footnotesize{}geophysical fluid dynamics}{\footnotesize{}
(3rd year)}{\footnotesize\par}
\end{itemize}

\column{9cm}

\includegraphics[width=8.3cm]{SEA2004/figures/heat-redistribution}
\begin{itemize}
\item \textbf{\footnotesize{}Heat is transported}{\footnotesize{} from the
equator to the poles through a network of atmospheric cells and ocean
currents}{\footnotesize\par}
\end{itemize}
\end{columns}

\end{frame}

\end{document}
